% cSpell:locale en,pt-BR - Adicionado para suporte ao corretor ortográfico
% Extensão Code Spell Checker no VS Code (muito útil!)
\chapter{Escrevendo o texto}\label{cap:texto}

Escolhido o tipo e preparado o ambiente, resta o texto. Esta seção trata das divisões do trabalho
--- e da única tradução mental que o modelo pede de você ---, das referências cruzadas entre elas e
das citações.

\section{O rebaixamento de níveis}\label{sec:rebaixamento}

\textbf{O comando que você digita não é o nome que o documento imprime.} É a única passagem do
modelo em que as duas coisas diferem, e vale conhecê-la antes de escrever a primeira divisão.

Este modelo é baseado na classe \texttt{book} do \LaTeX{}, cuja divisão de topo é o capítulo. A ABNT
NBR 14724:2024~\cite{nbr14724:2024}, no §4.2.2, não admite que um trabalho acadêmico seja dividido
em \emph{capítulos}: suas divisões chamam-se seções, da primária à quinária. O modelo concilia as
duas coisas rebaixando cada nível um degrau.

\begin{tabelaabnt}{tab:rebaixamento}{Comando digitado e nome impresso}
    \begin{tabular}{@{}llp{0.34\textwidth}@{}}
        \toprule
        \textbf{Você digita} & \textbf{O documento imprime} & \textbf{Prefixo de rótulo sugerido} \\
        \midrule
        \texttt{\textbackslash chapter\{\}}       & Seção (divisão primária) & \texttt{cap:} \\
        \texttt{\textbackslash section\{\}}       & Subseção                 & \texttt{sec:} \\
        \texttt{\textbackslash subsection\{\}}    & Subsubseção              & \texttt{ssc:} \\
        \texttt{\textbackslash subsubsection\{\}} & Parágrafo                & \texttt{sss:} \\
        \bottomrule
    \end{tabular}
    \fonte{o próprio autor}
\end{tabelaabnt}

A correspondência vale para os títulos, para o sumário e para as referências cruzadas. Ao escrever
``conforme a seção anterior'', lembre-se de que o comando que produziu aquela divisão foi o
\texttt{\textbackslash chapter}. Os prefixos de rótulo sugeridos são mnemônicos do \textbf{comando},
e não do nome impresso --- por isso \texttt{ssc:} marca aquilo que o documento chama de Subsubseção.

\subsection{Isto não é desvio da norma}\label{ssc:rebaixamento-nao-e-desvio}

O rebaixamento é \textbf{como} o modelo cumpre o §4.2.2, e não uma divergência dele. A saída ---
títulos, numeração, sumário --- é exatamente a que a norma descreve. O que muda é o vocabulário que
você digita, e ele é da classe \texttt{book}, não da ABNT.

A distinção importa porque este modelo \emph{tem} desvios deliberados, e eles são de outra natureza:
ali a saída difere do que a norma pede, por decisão registrada e justificada --- a folha de aprovação
sem linhas de assinatura é o mais visível deles (\refcomp{}{cap:estrutura}). Um desvio muda o que
sai no papel. O rebaixamento não muda nada no papel; muda só o que se digita.

\section{Estruturação de seções e subseções}\label{sec:secoes-desenvolvimento}

Quando se escreve um documento organizado, é importante estruturar o conteúdo em seções e subseções
de forma clara e lógica. Isso facilita a leitura e compreensão do texto, além de ajudar o leitor a
navegar pelo documento. Lembre-se que seu trabalho também servirá como referência para futuros
alunos e pesquisadores, portanto, uma boa estruturação é essencial.

As diretrizes de cada nível vêm a seguir. Os prefixos de rótulo sugeridos são os da
\reftable{tab:rebaixamento}, e são mnemônicos do comando digitado:


\begin{itemize}
    \item \textbf{Seções} (\verb|\chapter{}|): cada seção deve abordar um tema específico do trabalho.
          Lembre-se de que elas são numeradas automaticamente. Acrescente um rótulo logo após o título,
          com o comando \verb|\label{}|, para facilitar autorreferências dentro do documento. Por
          convenção, use o prefixo \texttt{cap:}, seguido de uma palavra-chave relacionada ao conteúdo
          (por exemplo, \texttt{cap:introducao} para a seção de introdução).
    \item \textbf{Subseções} (\verb|\section{}|): dentro de cada seção, dividem o conteúdo em partes
          menores e mais gerenciáveis. Também são numeradas automaticamente. Use o prefixo
          \texttt{sec:} para os rótulos (por exemplo, \texttt{sec:objetivos} para a subseção de
          objetivos).
    \item \textbf{Subsubseções} (\verb|\subsection{}|): caso seja necessário dividir ainda mais o
          conteúdo, organizam o texto em tópicos específicos e facilitam a leitura. Use o prefixo
          \texttt{ssc:} (por exemplo, \texttt{ssc:metodologia} para a subsubseção de metodologia).
    \item \textbf{Parágrafos} (\verb|\subsubsection{}|): a divisão mais detalhada, útil para tópicos
          que exigem explicação mais aprofundada. Use o prefixo \texttt{sss:} (por exemplo,
          \texttt{sss:detalhes-tecnicos}). É o nível mais baixo que o modelo numera, e seu uso é raro:
          detalhamento excessivo dificulta a leitura. Use com parcimônia.
    \item \textbf{Consistência}: mantenha uma estrutura consistente ao longo do documento. Use os
          mesmos níveis para temas semelhantes, garantindo que o leitor possa seguir o fluxo do texto
          facilmente.
\end{itemize}

Na sequência, veja como se parecem as subsubseções e os parágrafos.

\subsection{Subsubseção de exemplo}\label{ssc:exemplo}

Esta é uma subsubseção de exemplo. Note que ela é numerada automaticamente e pode ser
referenciada no texto usando o comando \verb|\ref{}| com o rótulo correspondente. Este modelo
também fornece o comando \verb|\refcomp{}{}|, que imprime o nome da divisão e seu título junto com o
número, no formato \emph{Divisão N - Título}. O nome da divisão é determinado automaticamente a
partir do \textbf{contador} do elemento referenciado, e não do primeiro argumento nem do texto do
rótulo --- ou seja, você não precisa (nem deve) digitá-lo, e renomear um rótulo não muda o nome
impresso. É por aqui que o rebaixamento de níveis se manifesta: um \verb|\subsection| referenciado
sai como \emph{Subsubseção}. Por exemplo, veja a seguir a referência para a divisão atual:
\refcomp{}{ssc:exemplo}. Note a formatação da saída e o \emph{hyperlink} criado automaticamente para
o elemento referenciado. O comando \verb|\ref{}| padrão do \LaTeX\ também funciona normalmente.

Existe ainda a forma com asterisco, \verb|\refcomp*{Nome}{rótulo}|, que imprime literalmente o nome
informado no primeiro argumento. Use-a apenas para rótulos fora da hierarquia acima --- contadores
próprios, ambientes personalizados etc. ---, para os quais o nome não pode ser determinado
automaticamente. Ela também serve quando você deseja um nome diferente do padrão: por exemplo,
\verb|\refcomp*{Ilustração}{fig:exemplo}| produz \refcomp*{Ilustração}{fig:exemplo}, em vez do nome
``Figura'' que seria usado normalmente.

\subsubsection{Parágrafo de exemplo}\label{sss:exemplo}

Assim como as demais divisões, os parágrafos também são numerados automaticamente e podem ser
referenciados no texto usando o comando \verb|\ref{}| com o rótulo correspondente. Por exemplo, veja
a seguir a referência para o parágrafo atual: \refcomp{}{sss:exemplo}.

\section{Citações}\label{sec:citacoes}

Nenhum trabalho acadêmico está realmente completo sem uma discussão adequada sobre as fontes que o
sustentam. Sendo assim, citações de trabalhos da literatura associada ao tema são fundamentais para
que o trabalho tenha credibilidade e rigor acadêmico. Duas normas da Associação Brasileira de Normas
Técnicas (ABNT) governam esse terreno, e convém não as confundir: a NBR
10520:2023~\cite{nbr10520:2023} trata da \textbf{citação no documento} --- a forma como o autor
remete, no corpo do texto, à obra que sustenta o que ele afirma ---, enquanto a NBR
6023:2018~\cite{nbr6023:2018} trata da \textbf{elaboração das referências} --- a lista, ao final do
trabalho, que descreve cada obra citada. Uma diz como citar; a outra, como descrever o que foi
citado. Este modelo utiliza o pacote \texttt{biblatex}%
%
\footnote{Este modelo descreve de forma muito sucinta o uso do pacote \texttt{biblatex}.
    Recomenda-se a consulta à documentação do pacote para completo entendimento de como utilizá-lo}%
%
 para gerenciar as referências, usando o pacote de estilo \texttt{abnt}, o que facilita a inclusão e
formatação das citações de acordo com essas normas. Sua bibliografia deve estar organizada em um
arquivo \texttt{.bib}, que neste modelo é denominado \texttt{referencias.bib} e está localizado na
raiz do projeto. \textbf{Atenção!}  O arquivo incluído é fictício, usando referências falsas apenas
para servir de exemplo de como organizar suas fontes.

Você também pode adicionar a bibliografia diretamente no arquivo principal do
documento, embora isso não seja recomendado para trabalhos mais extensos. Existem programas
gerenciadores de referências, como o \texttt{Zotero} e o \texttt{Mendeley}, que são úteis para
organizar suas fontes e são capazes de exportar suas referências diretamente em para um arquivo
\texttt{.bib}, facilitando a integração com o \LaTeX. Caso use uma dessas ferramentas, certifique-se
de revisar o arquivo exportado para garantir que todas as informações estejam corretas e completas e
simplesmente substitua o arquivo \texttt{referencias.bib} deste modelo pelo seu arquivo exportado.
Lembre-se de importá-lo no preâmbulo do arquivo principal do documento ou apenas o renomeie para
\texttt{referencias.bib}, substituindo o arquivo original. A seguir, alguns exemplos de como fazer citações corretamente no texto.

O primeiro tipo de citação é a citação direta, que consiste em transcrever exatamente as palavras de
um autor. Para isso, utilize o comando \texttt{\textbackslash cite} para incluir a referência no
texto. Citações diretas com até três linhas devem ser incorporadas ao parágrafo, entre aspas. Por
exemplo, segundo \textcite{fernandes2019}, ``a utilização de modelos padronizados em relatórios
técnicos é essencial para garantir a clareza e a uniformidade na apresentação dos resultados''. Note
que em \LaTeX, as aspas iniciais são representadas por dois acentos graves (\texttt{``}) e as aspas
finais por dois apóstrofos (\texttt{''}). Aspas simples são representadas por um acento grave
(\texttt{`}) e um apóstrofo (\texttt{'}) para abrir e fechar, respectivamente. Use aspas simples
para uma citação dentro de outra citação.

Citações diretas com mais de três linhas devem ser formatadas em um bloco separado. A ABNT NBR
10520:2023~\cite{nbr10520:2023} §7.1.1 pede quatro destaques ao mesmo tempo, e o \verb|\quote| dá os
quatro: recuo de 4 cm da margem esquerda, letra menor que a do corpo, \textbf{espaçamento simples}
--- e não o 1,5 do texto corrido, que a ABNT NBR 14724:2024~\cite{nbr14724:2024} §5.2 excetua
justamente aqui --- e ausência de aspas, que marcam apenas a citação curta. O bloco sai mais
compacto que o texto ao redor de propósito: não é erro de formatação a corrigir.
Supressões (quando você inicia a citação no meio de uma frase) devem ser indicadas por reticências entre colchetes
(\texttt{[\ldots]}) e acréscimos (quando você insere palavras para esclarecer o sentido da citação)
devem ser indicados por colchetes simples (\texttt{[texto adicional]}). Este modelo fornece o
comando \texttt{\textbackslash quote{}{}} para facilitar a formatação de citações longas. Veja o exemplo
a seguir:

\quote{martins2021}{%
[\ldots] modelos padronizados para elaboração de relatórios técnicos são fundamentais para a comunicação
científica eficaz. A estruturação comum garante que as ideias e fatos sejam devidamente comunicadas
aos receptores. A organização adequada do texto também é essencial para sua catalogação,
possibilitando a futuros leitores a compreenssão [grafia errada no original] do conteúdo de forma mais ágil e eficiente, assim
como sua consulta posterior.
}

Note que citações diretas longas não devem ser usadas com frequência excessiva, para evitar que o
texto se torne fragmentado e difícil de ler. Use-as apenas quando for realmente necessário
reproduzir fielmente as ideias do autor.

Sendo assim, o tipo mais comum de citação é a citação indireta,
que consiste em parafrasear as ideias de um autor com suas próprias palavras. Sintetize o conteúdo e
apresente-o de forma clara e concisa, sempre respeitando e atribuindo a fonte original.
Para fazer uma citação indireta, utilize o comando \texttt{\textbackslash textcite} ou
\texttt{\textbackslash parencite}, dependendo do contexto. Por exemplo, \textcite{silva2020} destaca a importância de
utilizar fontes confiáveis e atualizadas para embasar as argumentações em trabalhos acadêmicos.
Alternativamente, pode-se dizer que a utilização de fontes confiáveis é crucial para a credibilidade
do trabalho (cf. \parencite{silva2020}).

Existe ainda a \emph{citação de citação}, que ocorre quando você deseja citar uma fonte que foi
mencionada em outra fonte. Neste caso, usa-se a palavra ``apud'' para indicar que a citação foi
retirada de uma fonte secundária. No entanto, é importante ressaltar que a citação de citação deve
ser evitada sempre que possível, dando preferência à consulta direta da fonte original. Caso seja
necessário utilizar uma citação de citação, faça-o da seguinte forma: Segundo \textcite[p. 45]{rodrigues2017}, a padronização de relatórios técnicos é essencial para garantir a clareza e a uniformidade na apresentação dos resultados (apud \textcite{lima2022}).

\subsection{Exemplos de citações por tipo de referência}

A seguir apresentamos exemplos curtos, para cada tipo de referência incluída no arquivo
\texttt{referencias.bib}, mostrando uma citação indireta (paráfrase) e uma citação direta (trecho
transcrito) usando os comandos disponíveis neste modelo%
\footnote{Observação: os exemplos a seguir são fictícios e servem apenas para ilustração. As
    citações diretas usadas aqui tem menos de 3 linhas, são apenas exemplos. Lembre-se que você deve
    usá-las apenas quando a transcrição é superior à 3 linhas.}%
.

\begin{itemize}
    \item \textbf{Artigo (article)} --- Indireta: \textcite{article1} discute técnicas de amostragem que facilitam a inferência em redes complexas. Direta:\par
          \quote{article1}{A amostragem aleatória aplicada a topologias complexas reduz o viés nas estimativas de grau médio.}

    \item \textbf{Livro (book)} --- Indireta: Segundo \textcite{book1}, métodos estatísticos clássicos ainda são válidos em muitos problemas práticos. Direta:\par
          \quote{book1}{``Os testes de hipóteses constituem ferramenta central na análise de dados experimentais.''}

    \item \textbf{Anais / Trabalhos em Evento (inproceedings)} --- Indireta: \textcite{inproceedings1} apresenta otimizações aplicadas ao transporte urbano. Direta:\par
          \quote{inproceedings1}{``Aplicar heurísticas de baixo custo pode reduzir congestionamentos sem investimentos pesados.''}

    \item \textbf{Capítulo em livro (incollection)} --- Indireta: \textcite{incollection1} descreve modelos de previsão de demanda. Direta:\par
          \quote{incollection1}{``Modelos probabilísticos permitem estimativas robustas mesmo com dados esparsos.''}

    \item \textbf{Capítulo específico (inbook)} --- Indireta: \textcite{inbook1} cobre controle de qualidade na produção. Direta:\par
          \quote{inbook1}{``O controle estatístico de processo reduz a variabilidade de produção quando bem aplicado.''}

    \item \textbf{Tese/Doutorado (phdthesis)} --- Indireta: Em sua tese, \textcite{phdthesis1} explora métodos numéricos para dinâmica de fluidos. Direta:\par
          \quote{phdthesis1}{``A convergência do método implícito foi observada mesmo em malhas não-uniformes.''}

    \item \textbf{Dissertação/Mestrado (mastersthesis)} --- Indireta: \textcite{mastersthesis1} avalia algoritmos de roteamento. Direta:\par
          \quote{mastersthesis1}{``Algoritmos baseados em demand-aware routing apresentam melhor desempenho sob cargas variáveis.''}

    \item \textbf{Relatório Técnico (techreport)} --- Indireta: O relatório do \textcite{techreport1} sumariza aplicações de sensoriamento remoto. Direta:\par
          \quote{techreport1}{``Sensoriamento remoto é essencial para mapeamento ambiental em larga escala.''}

    \item \textbf{Manual (manual)} --- Indireta: O \textcite{manual1} descreve procedimentos de uso do sistema. Direta:\par
          \quote{manual1}{``Siga o procedimento de inicialização para evitar perda de dados durante a atualização.''}

    \item \textbf{Recurso online (online)} --- Indireta: Dados públicos do \textcite{online1} podem ser usados para análises temporais. Direta:\par
          \quote{online1}{``A série histórica de indicadores permite estudos comparativos entre regiões.''}

    \item \textbf{Material diverso (misc)} --- Indireta: \textcite{misc1} apresenta notas de aula úteis para introdução às redes neurais. Direta:\par
          \quote{misc1}{``A normalização das entradas melhora a estabilidade do treinamento.''}

    \item \textbf{Anais (proceedings)} --- Indireta: Os \textcite{proceedings1} reúnem trabalhos relevantes da área. Direta:\par
          \quote{proceedings1}{``Os anais deste simpósio apresentam avanços interdisciplinares em computação.''}

    \item \textbf{Patente (patent)} --- Indireta: A patente \textcite{patent1} descreve processo industrial inovador. Direta:\par
          \quote{patent1}{``O método aumenta a resistência do compósito sem aumentar o custo de produção.''}

    \item \textbf{Documento não publicado (unpublished)} --- Indireta: \textcite{unpublished1} traz resultados preliminares sobre microplásticos. Direta:\par
          \quote{unpublished1}{``Estudos iniciais indicam concentração elevada de partículas em amostras costeiras.''}
\end{itemize}
